\subsection{A mixed five-cycle family}

\begin{lemma}
If $n\geq3$ and $U\subseteq\mathbb C^n$ is linear with $e_{n-1}|_U=0$, then $\dim U\leq n-2$.
\end{lemma}
\begin{proof}
Let $\ell_1,\ldots,\ell_n$ be the coordinate linear forms on $U$.
If one is zero, at least two are zero: with exactly one zero, the sole surviving term of $e_{n-1}$ is the nonzero product of the remaining forms. Thus $\dim U\leq n-2$ in this case.
Otherwise partition the forms into proportionality classes, writing $\ell_i=a_iL_C$ for $i\in C$, where $a_i\ne0$. In the fraction field of $\mathbb C[U]$,
\[
 0=\frac{e_{n-1}(\ell)}{\ell_1\cdots\ell_n}
   =\sum_C\frac{\sum_{i\in C}a_i^{-1}}{L_C}.
\]
Multiply by $\prod_C L_C$ and reduce modulo each $L_C$.
Distinct representatives are relatively prime, so $\sum_{i\in C}a_i^{-1}=0$ for every class.
Each class has at least two members. Since the coordinate forms span $U^*$, we obtain $\dim U\leq\lfloor n/2\rfloor\leq n-2$.
\end{proof}

\begin{proposition}
Let $\omega$ be a primitive fifth root of unity. Every level space of six-by-six matrices whose rows have the form
\[
 c,\omega c,\omega^2c,\omega^3c,\omega^4c,d,
 \qquad c,d\in\mathbb C^6,
\]
has dimension at most nine.
\end{proposition}
\begin{proof}
The rank-one formula and $\prod_{a=0}^4(1+z\omega^a)=1+z^5$ give
\[
 S_1=\sum_jd_j,\qquad S_5=5!e_5(c).
\]
For the second identity, terms involving the last row and four of the first five rows have coefficient $e_4(1,\omega,\ldots,\omega^4)=0$.
The projection of the level space to $c$ therefore has dimension at most four by the lemma.
Its kernel consists of last-row vectors of sum zero and has dimension at most five. Rank-nullity proves the bound.
\end{proof}

\subsection{A small-characteristic obstruction to specialization certificates}

Consider the closed subscheme $Z_d\subseteq\operatorname{Gr}(d,36)_{\mathbb Z}$ defined by vanishing of the restrictions of the integer polynomials $S_1,\ldots,S_6$ to the universal subspace.
The morphism $Z_d\to\operatorname{Spec}\mathbb Z$ is projective, hence proper.
If $Z_d$ has a complex point, its image contains the generic point of $\operatorname{Spec}\mathbb Z$; its closed image is then all of $\operatorname{Spec}\mathbb Z$.
Thus emptiness over $\overline{\mathbb F}_p$ at any one prime would prove emptiness over $\mathbb C$.
For $d=11$, the primes two, three and five cannot be used in this way.

\begin{proposition}
Over each of $\mathbb F_2,\mathbb F_3,\mathbb F_5$, the equations $S_1=\cdots=S_6=0$ contain an eleven-dimensional linear subspace of six-by-six matrices.
\end{proposition}
\begin{proof}
In characteristic $p$, a block of $p$ identical rows contributes zero to every rook sum involving a nonempty selection from that block.
Indeed cyclic rotation of its row positions acts freely on the corresponding matchings and leaves each monomial value unchanged.
A nonempty matching cannot be fixed by a nonidentity rotation: invariance would assign the same column to all rows in a nontrivial row orbit, contradicting distinctness of columns.
Each orbit therefore contributes $p$ identical terms, hence zero.

In characteristic two, three independently chosen identical-row pairs give a space of dimension eighteen on which every positive rook sum vanishes.
In characteristic three, two independently chosen identical-row triples give dimension twelve.
In characteristic five, choose the first five rows equal and arbitrary, and choose the last row arbitrarily subject to its coordinate sum being zero. This gives dimension $6+5=11$.
The terms using the five-row block vanish, while the only possible nonempty matching confined to the last row contributes its zero first rook sum.
Taking an eleven-dimensional subspace in the first two constructions completes the proof.
\end{proof}

These are vanishing statements for the individual integer-defined layers. They do not use the equality $\operatorname{perm}(J+t)=6!$, which degenerates in these characteristics. An emptiness certificate by this specialization method must use a prime at least seven; no such certificate is proved here.

\subsection{Exploratory unrestricted numerical search}

An original numerical search was carried out on complex isotropic Grassmannian charts that impose $S_1=S_2=0$ up to floating-point roundoff. On the $35$-dimensional hyperplane $S_1=0$, the quadratic matrix is $Q=I-R^{\mathsf T}R-C^{\mathsf T}C$, where $R,C$ take row and column sums. After reducing this nondegenerate form to the identity, a Witt chart has basis
\[
 T=E+F\left(K-\tfrac12B^{\mathsf T}B\right)+HB,
 \qquad K^{\mathsf T}=-K.
\]
Here the columns of $E$ and $F$ are dual eleven-dimensional isotropic frames, and the columns of $H$ form an orthonormal frame for their thirteen-dimensional orthogonal complement. The matrix $K$ is eleven-by-eleven and skew-symmetric, and $B$ is thirteen-by-eleven. For dimension eleven, this gives $55+143=198$ complex parameters.
The implementation evaluates the permanent at seven roots of unity by Ryser's subset formula and extracts $S_3,\ldots,S_6$ by Fourier inversion, with analytic derivatives.
Each of sixteen starts used $1024$ training parameter vectors and an independent holdout of $512$ vectors after Hermitian orthonormalization of the final basis. Eight initial charts extended the known nine-dimensional hook by two isotropic directions, and eight used random charts; all chart parameters were allowed to vary.

No candidate was found. The smallest Euclidean norm of the four holdout root-mean-square layer residuals was $1.869$ for the hook-based starts and $4.206$ for the random starts. By comparison, a dimension-nine calibration with $512$ training points attained $1.65\times10^{-11}$. An earlier calibration with only $100$ training points had a tiny training residual but a holdout residual $4.14\times10^{-4}$; this detected overfitting is documented with the calculation.
The numerical kernel was independently compared with direct enumeration of all square subpermanents, with maximum absolute error $6.35\times10^{-12}$ on the check instances, and its analytic derivative was checked by centered finite differences.
The searches are discovery calculations, not infeasibility certificates: they neither cover the Grassmannian nor rule out other critical points or exact level spaces. Reproducible code, seeds, all start records and open CSV basis data accompany the paper.
